\section{The exact graph and comparison of actual returns}
\label{sec:graph}
The logarithmic form of the exact limiting return contracts the
complete independent state. We use it to construct an invariant graph
and a normalized scalar coordinate that records progress along its
returns. This coordinate fixes the explicit compliance sequences in
Theorem~\ref{thm:common-free}. The corresponding finite-scale contact and
free curves retain the order of motion needed for path comparison.

\begin{proposition}[An exact graph and normalized phase]
\label{prop:canonical-graph-phase}
After decreasing the outer radius $\delta_{\rm g}$, there are a common
radius $0<\delta_0<\delta_{\rm g}$ and a jointly
continuous graph $\Gamma_\lambda:(0,\delta_{\rm g}]\to\Kcal$,
with zero memory component, such that
\begin{equation}
 \begin{aligned}
 |\Gamma_\lambda(d)-\xi_{*,\lambda}|&\le Cd^{1/3},\\
 |\Gamma_\lambda(d_2)-\Gamma_\lambda(d_1)|
 &\le C\max(d_1,d_2)^{1/3}|\log(d_2/d_1)|.
 \end{aligned}\label{gr:graph-bounds}
\end{equation}
Its base map $F_\lambda(d)=f_{0,\lambda}(d,\Gamma_\lambda(d))$
is strictly increasing, $F_\lambda(d)>d$, and
$\sup_\lambda F_\lambda(\delta_0)<\delta_{\rm g}$. For
$0<d\le\delta_0$,
\begin{equation}
 T_{0,\lambda}(d,\Gamma_\lambda(d))
      =(F_\lambda(d),\Gamma_\lambda(F_\lambda(d))).
 \label{gr:invariance}
\end{equation}
On the outer interval there is a jointly continuous phase
$\mathcal S_\lambda$, strictly decreasing and bi-Lipschitz in $\log d$,
such that, with $\rho_c=3/4$,
\begin{equation}
 \mathcal S_\lambda(F_\lambda(d))
       =\rho_c\mathcal S_\lambda(d),\qquad
 \mathcal S_\lambda(d)=\log(K_\lambda^4/d)+O(d^{1/3}).
 \label{eq:canonical-Schroeder}
\end{equation}
The equation holds on the inner interval. The normalization fixes its
germ at zero on this graph. All bounds are uniform in $\lambda$.
\end{proposition}
\begin{proof}
Put $x=\log(K_\lambda^4/\delta)$ and
$\eta=\xi-\xi_{*,\lambda}$. Proposition~\ref{prop:limiting-factorization} gives
\begin{equation}
 \widehat T_\lambda(x,\eta)
  =(\rho_c x+r_\lambda(x,\eta),H_\lambda(x,\eta)),\qquad
 |r|+|H|+|Dr|+|DH|\le Ce^{-x/4}.
 \label{gr:log-map}
\end{equation}
For each $\lambda$ the domain
$[X,\infty)\times(\Omega-\xi_{*,\lambda})$ is convex. In the norm
$\|(x,\eta)\|=\max\{|x|,\|\eta\|_\infty\}$ its derivative row
sums are at most $\rho_c+Ce^{-X/4}$. Increase $X$ to obtain
\begin{equation}
 \|\widehat T_\lambda(z)-\widehat T_\lambda(\widetilde z)\|
       \le q\|z-\widetilde z\|,\qquad q=7/8.
 \label{gr:contraction}
\end{equation}
The base segment here is logarithmic, hence geometric interpolation
of its endpoint scales. It lies in the auxiliary domain of the exact
reference; actual returns will be compared by
\eqref{loc:value} at their physical endpoints.

We seek a curve $V_\lambda(S)=(S+u_\lambda(S),v_\lambda(S))$,
$S\ge S_0$, whose parameter is multiplied by $\rho_c$ under the return:
\[
 V_\lambda(\rho_c S)=\widehat T_\lambda(V_\lambda(S))
 \qquad(S\ge S_0/\rho_c).
\]
Thus $S$ will be the phase coordinate. To solve for the value at $S$,
evaluate the right-hand side at $S/\rho_c$. The error
$e^{-x/4}$ in \eqref{gr:log-map} then has size
$e^{-S/(4\rho_c)}=e^{-S/3}$, which determines the weight below.
On continuous pairs for $S\ge S_0$ set
$\|Z\|_0=\sup_{S\ge S_0}e^{S/3}(|u(S)|+\|v(S)\|_\infty)$,
and, writing $t=S/\rho_c$, define
\[
 (\mathcal PZ)(S)=
 \bigl(\rho_c u(t)+r_\lambda(t+u(t),v(t)),
                 H_\lambda(t+u(t),v(t))\bigr).
\]
On any fixed norm ball, all arguments lie in the common auxiliary
domain for large $S_0$. Since
$e^{S/3-t/4}=1$ and $e^{S/3-t/3}=e^{-S/9}$,
\[
 \|\mathcal P0\|_0\le C,\qquad
 \|\mathcal PZ-\mathcal P\widetilde Z\|_0
 \le(\rho_c e^{-S_0/9}+C_Me^{-4S_0/9})
                       \|Z-\widetilde Z\|_0.
\]
Choose $M\ge\max(1,2C)$ and then $S_0$ so the last coefficient
is below $1/2$. The fixed point satisfies
\begin{equation}
 u,v=O(e^{-S/3}),\qquad
 V_\lambda(S)=\widehat T_\lambda(V_\lambda(S/\rho_c)).
 \label{gr:parametrized-graph}
\end{equation}

The Picard iterates from zero are $C^1$. Their weighted derivative
norms $A_n=\sup e^{S/3}(|u_n'|+\|v_n'\|_\infty)$ obey
\[
 A_{n+1}\le C+(e^{-S_0/9}+C_Me^{-4S_0/9})A_n
\]
by the chain rule and \eqref{gr:log-map}. They are uniformly bounded.
Passing to the uniform limit in the integral inequality gives
\begin{equation}
 |u(b)-u(a)|+\|v(b)-v(a)\|_\infty
       \le C\int_a^b e^{-t/3}\dd t\quad(S_0\le a<b).
 \label{gr:weighted-Lipschitz}
\end{equation}
Finite iterates are jointly continuous in $(\lambda,S)$ on compact
intervals, and their common contraction tail preserves this continuity.
This argument uses only first state derivatives of the limiting return.

For larger $S_0$, choose $\eta<1/2$ with $Ce^{-S_0/3}\le\eta$.
Writing $x=S+u_\lambda(S)$, \eqref{gr:weighted-Lipschitz} gives
\begin{equation}
 (1-\eta)|S_2-S_1|\le|x_2-x_1|
                  \le(1+\eta)|S_2-S_1|.
 \label{gr:phase-bilip}
\end{equation}
Thus $S\mapsto x$ is strictly increasing and bi-Lipschitz.
On each fixed positive base interval this also gives two-sided
Lipschitz control between the phase and the base itself. Define
\begin{equation}
 \delta_\lambda(S)=K_\lambda^4e^{-S-u_\lambda(S)},\quad
 \Gamma_\lambda(\delta_\lambda(S))=\xi_{*,\lambda}+v_\lambda(S),\quad
 \mathcal S_\lambda(\delta_\lambda(S))=S.
 \label{gr:phase-definition}
\end{equation}
Choose $\delta_{\rm g}<\min_\lambda \delta_\lambda(S_0)$ and
$\delta_0<\min_\lambda \delta_\lambda(S_0/\rho_c)$, within the
limiting-return radii. Reduce $\delta_0$ so that
$\sup_\lambda F_\lambda(\delta_0)<\delta_{\rm g}$.
Equation \eqref{gr:parametrized-graph} gives
$F_\lambda(\delta_\lambda(S))=\delta_\lambda(\rho_c S)>\delta_\lambda(S)$
on the inner interval. Both factors are strictly decreasing in their
arguments, so $F_\lambda$ is strictly increasing there. It remains
strictly increasing on the outer interval after reducing its radius:
\eqref{loc:factorization} and \eqref{gr:weighted-Lipschitz} give a
positive lower bound for its logarithmic difference quotient.
The same estimates give \eqref{gr:graph-bounds} and \eqref{eq:canonical-Schroeder}.
The memory component of the fibre map vanishes, hence so does
that of the graph.
For two coordinates with the stated normalization on this graph,
their difference satisfies
$E(d)=\rho_c^jE(F_\lambda^{-j}(d))\to0$, proving uniqueness
of the normalized germ. The unconstrained part of the outer interval
is only a representative of that germ.
\end{proof}

The graph's logarithmic Lipschitz constant on $(0,\delta_{\rm g}]$
is $L_\Gamma\le C\delta_{\rm g}^{1/3}$. Reduce this outer radius,
and then $\delta_0$, so that $q(1+L_\Gamma)<1$. With the common
graph domains fixed, choose $\delta_{\rm reset}$ last as in
Theorem~\ref{thm:detuned-output-admission}, and take $s_0>0$ with
$2s_0<\min(\delta_0,\delta_{\rm reset})$. All reference curves with
input in a fixed positive compact subset of $(0,2s]$ are then
available in the same construction.

\subsection{The ordered one-return curve}
\label{gr:ordered-reference}

We spell out the finite-scale curve used in the path comparison.
Let $p(\phi_d)=-ad$ on the selected branch and put
$h_d=\mu A(\phi_d)-B(\phi_d)>0$. Write
$\Gamma_\lambda(d)=(0,\Gamma_U(d),\Gamma_W(d))$ and define
\begin{equation}
 \begin{gathered}
 U_d=d_\lambda d\{1+d^{1/3}\Gamma_U(d)\},\qquad
 V_d^{\rm in}=-r_0d^2\{1+d^{1/3}\Gamma_W(d)\},\\
 V_d^{\rm a}=V_d^{\rm in}-(ad/h_d)U_d,\qquad
 \mathcal E_\lambda(d,u,v)
     =(\phi_d,0,(u-U_*)/\cos\phi_d,v,0).
 \end{gathered}\label{gr:embedding}
\end{equation}
Here $u$ denotes the tuned slip argument of the embedding. Run the exact
reference compression at $\sigma=\sqrt d$:
\begin{equation}
 Y'=V_c,\quad \sigma Z'=Y-2Z,\quad
 V_c'=-a(Z-Y),\quad \mathcal U'=-h_d(Z-Y),
 \label{gr:compression}
\end{equation}
from $(Y,Z,V_c,\mathcal U)=(0,0,V_d^{\rm in}/d^2,U_d/d)$.
At its first arrest, start the inherited reference
$N_d'=A_DN_d$, $N_d=(Y_d,V_d^{\rm n},Z_d)$, with
$N_d(0)=(Y,\sigma V_c,Z)$ at arrest. Let $R_d$ be its first
force zero. The first-event margins from Section~\ref{sec:contact} give
\begin{equation}
 Z_d(R_d)=Y_d(R_d)<0,\quad V_d^{\rm n}(R_d)>0,\quad
 V_d^{\rm n}(R_d)+Z_d(R_d)>0.
 \label{gr:release-signs}
\end{equation}
Put $V_d^{\rm rel}=d^{3/2}V_d^{\rm n}(R_d)$ and
$\mathcal R_d=\mathcal E_\lambda(d,0,V_d^{\rm rel})$.
Write $\Phi_{\rm f}(\ell;\mathcal R_d)$ for the analytic free
mechanical flow, with zero memory appended. Its first subsequent
touchdown has elapsed time $T_d$ and scale $F_\lambda(d)$.

The continuous ordered curve is
\begin{equation}
 \cC_{\lambda,d}(\alpha)=
 \begin{cases}
 \mathcal E_\lambda(d,(1-\alpha)U_d,
                  V_d^{\rm in}-(ad/h_d)\alpha U_d),&0\le\alpha\le1,\\
 \mathcal E_\lambda(d,0,d^{3/2}V_d^{\rm n}((\alpha-1)R_d)),&1\le\alpha\le2,\\
 \Phi_{\rm f}((\alpha-2)T_d;\mathcal R_d),&2\le\alpha\le3.
 \end{cases}\label{eq:ps2-reference-path}
\end{equation}
The compression invariant
$V_c-(a/h_d)\mathcal U=V_d^{\rm in}/d^2-(a/h_d)U_d/d$
gives $d^{3/2}V_d^{\rm n}(0)=V_d^{\rm a}$, so the first join agrees.
The second join is $\mathcal R_d$. For $0<d\le\delta_0$, graph invariance gives
\begin{equation}
 \cC_{\lambda,d}(3)=\cC_{\lambda,F_\lambda(d)}(0),\qquad
 \sup_\alpha|\cC_{\lambda,d}(\alpha)-G|\le Cd^{3/4}.
 \label{gr:curve-joins}
\end{equation}
The contact pieces have instantaneous physical duration in this
reference curve; their parameters preserve the order of compression
and release. They are not extra free solutions. At every fixed
positive compact input interval, including the enlarged intervals
used to scan $[s/2,2s]$, the event maps, free flow and graph bound
make \eqref{eq:ps2-reference-path} Lipschitz in its input, uniformly in $\lambda$.
The corresponding constants may depend on that interval.

The same radius choice makes $U_d>0$. Consequently the contact
curve is injective on $[0,2]$: the tuned slip $(1-\alpha)U_d$
strictly decreases on compression, while on release it is zero
and the normal velocity satisfies
\begin{equation}
 \frac{\dd V_d^{\rm n}}{\dd\tau}
       =\tfrac12(Z_d-Y_d)>0\qquad(0\le\tau<R_d).
 \label{gr:contact-injectivity}
\end{equation}
Positivity of the active force gives the strict inequality, and
the two pieces meet only at their common arrest point. We will use
this fact both to identify an interior stopping target and to
separate complete path images.

\subsection{The free-section parameter}
\label{gr:free-section}

Along the free part, $\dot p=-a+O(d+\ell)<-a/2$ after the common
radius is reduced. Thus $p$ strictly decreases from $-ad$ to
$-aF_\lambda(d)$. Fix $0<s\le s_0$ and put
$d_*=F_\lambda^{-1}(s)$; this inverse is defined since $F$ is
continuous and strictly increasing from zero and $F(s)>s$.
For $d_*\le d\le s$, the equation
\begin{equation}
 p(\Phi_{\rm f}(\ell_s(d);\mathcal R_d))=-as
 \label{gr:section-root}
\end{equation}
has exactly one root in $[0,T_d]$. Its derivative in $\ell$ is
bounded away from zero, so its continuous dependence, including
the endpoints, follows either from the inverse function theorem
or directly from strict monotonicity and compactness. Hence
\begin{equation}
 j_s(d)=2+\ell_s(d)/T_d,
 \qquad j_s(d_*)=3,\qquad j_s(s)=2
 \label{gr:free-section-definition}
\end{equation}
is the continuous parameter in \eqref{eq:path-free-section}.
On contact the section value is instead the constant $p=-ad$;
therefore its section points have $d=s$ and $0\le\alpha\le2$.
Together these observations give exactly the terminal domain
$\mathcal A_s$ of \eqref{eq:ps5-domain}.
